\addchap{Preface}
\begin{refsection}

\textit{Topics in the semantics of Slavic languages} contains twelve chapters devoted to a particular topic in formal semantics, as well as an introductory chapter. We are also grateful to Hana Filip, a key figure in the history of Slavic formal semantics, for contributing a personal note, which follows this preface. The topic chapters all consist of two parts: an overview part and a new analytical contribution or case study. The volume fulfills a number of goals: (i) to summarize existing analyses and theories, (ii) to provide a representative bibliographic overview of the relevant literature, (iii) to demonstrate what Slavic data and theories built upon them have to offer to the general linguistic discourse, and (iv) to provide a novel theoretical and/or empirical contribution. 

The collection is of interest to a broad readership -- Slavic as well as general linguistic scholars, experienced researchers as well as advanced undergraduate or postgraduate students. Individual contributions can serve as background reading to topical seminars. The volume covers areas of semantics in which Slavic languages have traditionally or more recently played an important role in the general formal semantic discourse (e.g. aspect, event structure, genericity, tense, information structure, negation and NPIs, (in)definiteness and specificity, bare NP semantics), as well as topics that have received less attention in their application to Slavic languages, but are pertinent to the theory of semantics, including its interfaces (e.g. focus participles, questions, imperatives, numerals, passives).   

All the chapters have solid empirical grounding and, where applicable, the authors emphasize the existing empirical generalizations. The use of state-of-the-art empirical methods, including corpus and experimental research, were encouraged, but not required. Although the focus of the volume lies on semantics, authors were free (and in some cases encouraged) to discuss the pertinent semantic issues in the context of any relevant interface, especially morphosyntax and pragmatics. Authors were encouraged to include cross-Slavic comparisons and to include discussion of Slavic languages that are less well studied. The chapters in this collection discuss data from the following Slavic languages: Bosnian/Croatian/Montenegrin/Serbian (BCMS), and more specifically Croatian and Serbian, Bulgarian, Czech, Macedonian, Polish, Resian, Russian, Slovak, Slovenian, Upper Sorbian, and Ukrainian. The contributions successfully completed a one-round review process in which every article was evaluated and commented on by one internal reviewer, i.e. someone who also contributed a chapter, as well as two external reviewers, one of whom was a PhD student at the time of reviewing.

This book would not have been possible without the help of the following reviewers: Daniel Altshuler, Maša Bešlin, Petr Biskup, Joanna Błaszczak, Olga Borik, Donka Farkas, Hana Filip, Berit Gehrke, Ljudmila Geist, Nina Haslinger, Tania Ionin, Peter Jenks, Dorota Klimek-Jankowska, Natasha Korotkova, Arkadiusz Kwapiszewski, Dimitra Lazaridou-Chatzigoga, Fabienne Martin, Ora Matushansky, Sergey Minor, Olav Mueller-Reichau, Rick Nouwen, Stefan Milosavljević, Maria Onoeva, Mareike Philipp, 
Agata Renans, Deniz Rudin, Radek Šimík, Frank Staniszewski, Peter R. Sutton, Elena Vaikšnoraitė, Marcin Wągiel, Hedde Zeijlstra, and Karolina Zuchewicz. Many thanks for your reviews! Furthermore, we thank the following people for assisting in \LaTeX {} typesetting and not leaving all that work to us two: Olga Borik, Mojmír Dočekal, Hana Filip, Atle Grønn, Olav Mueller-Reichau, Adrian Stegovec, and Marcin Wągiel. We also wish to acknowledge the technical support of the Language Science Press editorial team, especially Sebastian Nordhoff, Felix Kopecky, as well as Matthew Korte for proofreading the whole book. Finally, Radek Šimík would like to acknowledge the institutional funding Cooperatio LING (Faculty of Arts, Charles University) for enabling the successful completion of this book project.%the help of those who contributed by %type-setting and RADEK: TYPE-SETTERS HAVE A DEDICATED PLACE WHERE THEY'RE MENTIONED, I THINK
%proofreading parts of the contents of this book. 

We initiated this project in 2019 with the aim of having the volume published earlier than it was published now, but unforeseeable circumstances (including the pandemic) led to delays, both on the side of the authors and editors. Nevertheless, we hope that the readers will find it interesting and inspiring.\\

\hfill Berit Gehrke

\hfill Radek \v{S}imík

\hfill Berlin \& Prague, March 2026


% {\sloppy\printbibliography[heading=subbibliography]}
\end{refsection}

